% show abstract on the title page
\begin{abstract}

Population ageing has raised concerns about the future of aggregate labour supply, but its effects depend not only on the rising share of older individuals but also on immigration, participation, hours worked, and the changing age composition of the working-age population. This article estimates aggregate labour supply in Spain between 2005 and 2025 using Labour Force Survey microdata and decomposes its change by age group and birthplace status. 

\vspace{0.7em}\par
The results show that Spain's labour supply shifted strongly away from young workers and toward prime-age and older workers. The contribution of young workers aged 15--34 declined more than their population share, indicating weaker labour-supply attachment in addition to demographic contraction. By contrast, older workers aged 55+ increased their labour-supply share more than their population share, suggesting stronger labour-market attachment among older individuals. Immigration played a decisive compensating role: foreign prime-age and older workers generated substantial labour-supply growth, while the younger age structure of foreigners offset a large share of the negative age-structure effect associated with the ageing of the Spanish population. Overall, the findings show that ageing does not mechanically reduce aggregate labour supply; its effect is mediated by immigration, participation behaviour, and hours worked.

\vspace{0.7em}\par
\textbf{Keywords}: Labour Supply, Immigration, Population Ageing, Decomposition Models.
\end{abstract}

